\subsection{Fourier cancellation in the nonlinear remainder}
\label{app:enclosing-domain}
The higher-band failure motivates a further exploratory refinement of the
\emph{same} continuous-density class. The preceding frequency-wise Taylor
bound discards cancellation in the known adjoint field. For one particle let
\(f(x)=\operatorname{Re}\sum_q c_q\exp(2\pi i k_q^T x)\), where
\(c_q=(w_{Rq}+i w_{Iq})C_q/\sigma\). Let \(R=\sqrt{3}/2\), let \(a\) be
the rotation radius, and let \(s_i\) bound the norm of the three-dimensional
translation dual to the detector shift. For exact isometric detector
embeddings, \(s_i=s\); in computation we use the embedding pseudoinverse and
an explicit residual-phase correction, as recorded with the implementation.

For a joint scaled pose ball, the path \(y(t)=\exp(t\Omega)x+tb\),
\(0\leq t\leq1\), satisfies
\begin{equation}
 L=\sqrt{a^2R^2+s_i^2},\qquad H=a^2R,\qquad T=a^3R
\end{equation}
as bounds on its first, second and third derivative norms. An enclosing ball
has radius \(R+s_i\). An alternative enclosing coordinate cube has half-side
\begin{equation}
 r_i=\tfrac12+2R\sin\{\min(a,\pi)/2\}+s_i.
\end{equation}
Indeed, \(\norm{(\exp(t\Omega)-I)x}\leq
2R\sin\{\min(a,\pi)/2\}\), and translation adds at most \(s_i\).
This enlarges the integration domain of a \emph{known} field; it imposes no
new support or smoothness condition on the unknown density.

\begin{proposition}[Enclosing-domain remainder]
Let \(D_i\) contain every transformed unit cube along every allowed path,
and let \(S_m=\norm{D^m f}_{L^2(D_i),F}\) be the Frobenius derivative-tensor
norm. The second-order pose Taylor remainder has \(L^2\) norm at most
\begin{equation}
 r_{\mathrm{enc},i}=\tfrac16\{L^3 S_3+3LH S_2+T S_1\}.
 \label{eq:enclosing-remainder}
\end{equation}
\end{proposition}
\begin{proof}
The third chain-rule derivative is the sum of the cubic derivative contraction,
three Hessian contractions with \(y'\) and \(y''\), and the gradient
contraction with \(y'''\). Cauchy--Schwarz bounds each contraction by its
Frobenius norm. The map from \(x\) to \(y(t)\) is rigid and has Jacobian one;
enlarging its image to \(D_i\) increases each nonnegative squared integral.
Minkowski's inequality and Taylor's integral remainder give
\eqref{eq:enclosing-remainder}.
\end{proof}

For either symmetric domain, write \(K_D(v)=\int_D e^{2\pi i v^Tx}\,dx\).
The exact real-arithmetic identity retaining cross-frequency terms is
\begin{align}
 S_m^2={}&\tfrac12(2\pi)^{2m}\operatorname{Re}
 \sum_{q,l}(k_q^Tk_l)^m\bigl[
 c_q\overline{c_l}K_D(k_q-k_l)\nonumber\\
 &\hspace{32mm}+(-1)^m c_qc_l K_D(k_q+k_l)\bigr].
\end{align}
The sum over ordered tensor derivatives yields \((k_q^Tk_l)^m\);
the sum-frequency sign is \(i^{2m}=(-1)^m\). For a ball of radius \(r\),
\(K_D(v)=4\pi r^3j_1(z)/z\), \(z=2\pi r\norm{v}\), with value
\(4\pi r^3/3\) at zero. For the coordinate cube,
\(K_D(v)=\prod_j 2r\operatorname{sinc}(2rv_j)\).
The scalar density bias receives \((B+P)\sum_i r_{\mathrm{enc},i}\).
Taking the smaller valid remainder for each particle is deterministic and
uses the original spectral event; no new random certificate is selected.
Magnitude-based roundoff guards remain heuristic, not validated arithmetic.

Independent spherical and tensor-product quadrature checks all ordered
derivative tensors through degree three. Direct nonlinear cube integration
also checks allowed joint poses and deliberately nonplanar frequency
embeddings. On the saved 1,024-particle higher-band example, the ball enclosure
reduces the cubic contribution from 56.493 to 32.623 after particle-wise
selection, still yielding no data. The subsequent expanded-cube probe reduces
it to 21.110, with total half-width 25.840, or 0.751 of no data. All particles
improve in that probe; its extra embedding-residual bias is below
\(5\times10^{-14}\). Both probes take approximately five CPU seconds and
retain the previous polynomial certificate and estimator weights. Nevertheless,
the three reference scenarios still have numerically zero feature-detection
power. This inexpensive refinement reduces a conservative remainder; it does
not solve the fine-feature usefulness or experimental calibration problem.
